\documentclass[12pt]{article}
\usepackage{amsmath,amssymb}

\begin{document}
\section*{{2025 AMC 12A Problems/Problem 3}}
Source: AoPS Wiki

\subsection*{Solution}
When Ash joins a team, the team's average is pulled towards his age. Let $A$ be Ash's age and $N$ be the number of people on the student team. This means that there are $15-N$ people in the teacher team. Let us write an expression for the change in the average for each team. The students originally had an average of $12$ , which became $14$ when Ash joined, so there was an increase of $2$ . The term $A-12$ represents how much older Ash is compared to the average of the students'. If we divide this by $N+1$ , which is the number of people on the student team when Ash joins, we get the average change per team member once Ash is added. Therefore, \[\frac{A-12}{N+1} = 2.\] Similarly, for teachers, the average was originally $55$ , which decreased by $3$ to become $52$ when Ash joined. Intuitively, $55-A$ represents how much younger Ash is than the average age of the teachers. Dividing this by the expression $(15-N)+1$ , which is the new total number of people on the teacher team, represents the average change per team member once Ash joins. We can write the equation \[\frac{55-A}{16-N} = 3.\] To solve the system, multiply equation (1) by $N+1$ , and similarly multiply equation (2) by $16-N$ . Then add the equations together, canceling $A$ , leaving equation $43=50-N$ . From this we get $N=7$ and $A= \boxed{28}.$ ~lprado

\end{document}
